\section{Method: boundary, unit, table, class}
\label{sec:method}

A thread that claims to improve an attack must express its progress as
a ratio to a boundary, reported in a single table whose rows are
variants and whose columns are one unit. A thread that cannot state
its boundary has not started; a thread whose ratio to the boundary is
flat has not progressed, however much its constants moved. The rule is
the reporting convention of the sibling study
library.\art{crypto:AGENTS.md}

\subsection{Two boundaries}

A \emph{boundary} is a quantity the method under study cannot cross,
derived rather than measured. A thread normally has one of each kind.

\begin{itemize}[leftmargin=1.2em,itemsep=0.3em]
\item A \emph{floor}: a lower bound from a counting or generic-group
  argument. For residual walks on $E(\FF_{p^3})$ a generic algorithm's
  expected relation yield from $P$ residuals is at most $\gamma P^2/2n$,
  so the count factor obeys $\kappa_{\mathrm{total}}\ge\sqrt{2(B+1)/\gamma}$.
  That floor moves only with the factor-base size $B$ and the
  automorphism order $\gamma$; it cannot be tuned
  away.\art{crypto:research/notes/index-calculus/RESEARCH_RESIDUAL_WALKS.md, §9.6}
\item A \emph{reference}: the cost of the best algorithm that already
  solves the same problem, measured in the same unit on the same
  instances. For the discrete logarithm that is Pollard rho, run on
  the same group, with the same operation accounting.
\end{itemize}

Both are written down \emph{before} optimising. A boundary chosen after
the fact is not a boundary.

\subsection{One table, one unit}

The unit used throughout is
\[
  S \;=\; \frac{\text{total operations}}{\sqrt{n}}.
\]
Rho is then a flat $S\approx 1.3$ at every size, and ``is this better
than rho'' is one column. Every cost the method incurs belongs in $S$:
precomputation, table setup, relation verification, linear algebra,
and any work an oracle does per call. Foreign units are converted with
a measured factor (for instance $63$ field multiplications per curve
addition on $\FF_{p^3}$) and the factor is recorded.

End-to-end speed is the measure of speed. $S$ is defined over the
whole method, cold, from setup to the recovered logarithm. A number
that prices only one phase --- the relation search, the decomposition
oracle, a single solver call, the linear algebra --- is a \emph{stage
diagnostic}, never a speed. ``Faster than rho'' means the method's $S$
column, with every phase inside it, sits below rho's, robustly, at
growing $n$. The only admissible speed ratio is
\[
  \mathrm{speedup}
  \;=\;
  \frac{\text{baseline total operations}}
       {\text{candidate total operations}}.
\]
A ratio taken over any smaller slice of the pipeline is labelled a
stage diagnostic and is not reported as a speedup. A row without a
verified answer is not a result.

\subsection{Four classes}

Every change is classified by what moved, not by how it felt to make.

\begin{center}
\small
\begin{tabular}{@{}llp{0.46\textwidth}@{}}
\toprule
class & what moved & what it means \\
\midrule
\textbf{advance} & ratio to the floor fell
  & the method does something a generic algorithm cannot \\
\textbf{engineering} & $S$ fell, ratio to the floor flat
  & legitimate, bounded, and not a finding \\
\textbf{relabelling} & headline count fell, $S$ rose
  & work moved somewhere the headline does not look \\
\textbf{accounting} & numbers changed, algorithm did not
  & a correction; say so and do not claim the gain \\
\bottomrule
\end{tabular}
\end{center}

All four are worth committing. Only the first is a result.

\subsection{Worked case: residual walks}
\label{sec:method-residual}

The residual-walk note is the reference implementation of the rule.
Section~9.6 of that note states the boundary and the numeric target
($\kappa_{\mathrm{total}}/\kappa_{\mathrm{floor}}<0.9$ on a hybrid at
fixed $B$, correct answer on every seed, zero relations failing
verification, zero trivial collisions counted as relations).
Section~9.7 freezes the baseline table. Section~10.6 classifies every
lever against the floor. Sections~11.5--11.6 are an engineering ledger
with the class of each step named. Section~11.7 prices the phase the
earlier rounds had skipped and corrects the conclusion they
implied.\art{crypto:research/notes/index-calculus/RESEARCH_RESIDUAL_WALKS.md}

The relabelling class is not hypothetical. In §10.5 a
triple-decomposition oracle cut the walked count from $16.3$ to
$0.15$, a hundredfold move in the number the thread had been tracking,
while total cost got $4.6\times$ worse: the count was not reduced, it
was paid for per residual instead of per step.

The skipped-phase mistake is the other half. Three rounds priced the
relation phase, quoted a crossover with rho at $2^{96}$, and left the
linear algebra unpriced. When it was measured it grew as $n^{0.68}$
against relations at $n^{1/3}$, so the method's cost bottoms out near
$200\times$ rho around $2^{50}$ and rises after. The $2^{96}$ was a
statement about one phase and had been read as a statement about the
method. That is the §5 mistake of the rule, named so that this paper
does not repeat it.

On $E(\FF_{p^3})$ at the sizes that fit, rho costs $S\approx 1.3$ and
every index-calculus variant built costs between $528\times$ and
$4{,}000\times$ that, with the asymptotically better variant crossing
rho only past $2^{230}$ (extrapolation, exponents in
Section~\ref{sec:ic-cubic}).

\subsection{Worked case: the ECC2K-130 pipe}
\label{sec:method-pipe}

The second worked case is a throughput floor rather than an operation
count. Any ECC2K-130 walk that keeps the rho collision structure
performs one affine addition per iteration. Affine addition with
batched inversion costs a fixed number of field products; a $131\times
131$ carry-less product is six \texttt{CLMAD}. The carry-less unit on
the RTX~PRO~6000 issues at about $1.6$ lane-operations per SM-clock.
Those two facts write a floor in billions of iterations per second
before any kernel is tuned: $28$~B/s needs $\le 25.4$--$26.9$
\texttt{CLMAD} per update including squarings and inversion, and the
arithmetic floor is already $28.9$ even with squarings moved
elsewhere.\art{crypto:ecc2k130/ITERATION-FUNCTION.md, §1}

Every subsequent GPU row is engineering against that floor. The table
walk of Section~\ref{sec:tablewalk} declared $17.0$~B/s as its own
target before the build; it measured $16.56$ and was declined. That is
the rule working in the other direction: a pre-declared numeric
condition that a later run either meets or does not.

\subsection{What does not count}

A ratio improved by moving the boundary, changing the unit, or
dropping a cost from the budget; a count lowered without the total
falling; an extrapolation presented as a measurement; a run whose
answer was not checked against the planted secret; wall-clock time as
the headline. Extrapolate freely, and mark it as extrapolation with
the exponents it rests on. Fit exponents over at least four sizes and
report the fit next to the prediction.
